\section{Deviations from the analysis plan}
\label{app:cls_dev}
The deviations log (supplementary) records every change after the plan was frozen, with its date and reason. In short:
\begin{itemize}
\item \emph{Gradient accumulation.} Swin-V2-B did not fit a 16\,GB GPU at batch size 32 and was screened with four micro-batches per step
  (same effective batch). The other encoders ran unchanged.
\item \emph{Encoder learning rate in the screen.} ConvNeXt-V2-T stayed at chance at the plan's learning rate. The three encoders without
  batch normalisation were screened again with the encoder learning rate reduced fourfold; this was decided before the other screen results
  were seen, on validation data only.
\item \emph{Holm families.} Before any test evaluation the authors fixed how H6 and H7 would be corrected if two encoders went forward, and
  which models H10 uses. With one encoder the families are those of Section~\ref{sec:cls_stats}.
\item \emph{Execution.} Data loading, reruns of unfinished jobs and resumption across sessions changed without changing any analysis.
\item \emph{Repeated NIH runs.} One account's NIH soft-tissue array differed at the byte level from the other five copies. The five runs
  trained on it (soft seeds 0, 3, 4; early seeds 2, 3) were repeated on a consistent copy before the analysis. With the superseded runs,
  the early and soft seed ensembles would have scored 0.8279 and 0.8253 instead of 0.8273 and 0.8243; this check lies outside the plan and
  entered no test.
\end{itemize}

\section{Additional figures}
\begin{figure}[htbp]
\centering
\safefig{figures/fig_seed_spread_densenet121.pdf}
\caption{Test macro-AUROC of every run (open circles; BDCXR-3257: each run on its own test fold) and of the seed ensembles that the tests
compare (diamonds).}
\label{fig:cls_seeds}
\end{figure}

\begin{figure}[htbp]
\centering
\safefig{figures/fig_conformal_sets_densenet121.pdf}
\caption{Conformal prediction at nominal 0.90: (a) coverage per true class and (b) mean set size of the Mondrian sets on BDCXR-3257;
(c) test false-negative rate and (d) mean set size under false-negative risk control on NIH ChestX-ray14. Dashed: the nominal level.}
\label{fig:cls_conformal}
\end{figure}

\begin{figure}[htbp]
\centering
\safefig{figures/fig_xai_ig_densenet121.pdf}
\caption{Integrated Gradients per input stream of one hybrid model (seed 0) on the NIH ChestX-ray14 test images of
Fig.~\ref{fig:cls_maps}.}
\label{fig:cls_ig}
\end{figure}
